\chapter{O que ainda não sabemos}
\label{cap:futuro}

\begin{objetivos}
Ao final deste capítulo, você será capaz de:
\begin{itemize}[leftmargin=*,itemsep=0.2em]
  \item enunciar com precisão as limitações que atravessam todo o trabalho apresentado;
  \item distinguir limitações de dados, de método e de avaliação;
  \item formular questões de pesquisa a partir de resultados negativos e de artefatos
        identificados;
  \item avaliar criticamente afirmações sobre sistemas de detecção de anomalias em
        contextos industriais;
  \item planejar uma agenda de trabalho que ataque os gargalos na ordem correta.
\end{itemize}
\end{objetivos}

\section{Uma advertência sobre o que foi medido}

Este livro apresentou muitos números. Antes de propor o que vem depois, é necessário
delimitar o que esses números sustentam.

\begin{destaque}[A ressalva que atravessa tudo]
Quatro classes receberam, do agente, nomes contendo ``perda de telemetria''
(\cref{cap:resultados-agente}). Isso revelou que, no 3W, o início de muitas anomalias
coincide com sensores de fundo reportando valores ausentes ou zerados.

A implicação é desconfortável: \textbf{parte do desempenho reportado nos experimentos
com o 3W pode decorrer da detecção de falha de instrumentação, e não da detecção da
anomalia}. Não há como quantificar essa fração com os experimentos realizados.

As análises por matriz de confusão apresentadas não evidenciaram esse padrão; as
descrições em linguagem natural o tornaram visível e motivam um teste controlado.
\end{destaque}

Essa ressalva não invalida automaticamente os resultados --- o \cref{cap:dual} reporta
estudos de caso de campo com dados reais e eventos verificados --- mas define a primeira
prioridade da agenda: \emph{medir o confundimento antes de perseguir mais acurácia}.

\section{Limitações de dados}

\begin{enumerate}[leftmargin=*,itemsep=0.4em]

\item \textbf{Possível confundimento de telemetria.} Descrito acima. Testável por um protocolo
que remova instâncias em que a anomalia coincide com valores ausentes, e repita toda a
avaliação. A diferença de desempenho pareada, com intervalo de confiança, estimará a
contribuição desse sinal; a remoção também pode alterar a composição e a dificuldade da
amostra.

\item \textbf{Um único conjunto de dados.} Quase tudo foi avaliado no 3W. A exceção ---
o sistema dual, avaliado em poços reais de um operador --- é justamente a que não permite
publicação dos dados.

\item \textbf{Amostras pequenas por classe.} A classe 9 tem $n = 11$ no estudo de
novidade e $n = 14$ no de classificação. Intervalos de confiança de Wilson com essas
amostras são largos o bastante para tornar várias comparações inconclusivas.

\item \textbf{Rótulos de qualidade heterogênea.} A análise do poço A3 no
\cref{cap:segmentacao} mostrou que a discordância entre modelo e anotação podia ser
atribuída à anotação. Não se sabe quantos outros casos assim existem.

\item \textbf{Classes definidas por sintoma.} As classes 4 e 5, discutidas em quatro
capítulos. Não é um problema de método.

\item \textbf{Possível contaminação.} O 3W é público desde 2019 e pode ter integrado
corpora de pré-treinamento. O \cref{cap:resultados-agente} apresentou três argumentos
contra essa explicação, nenhum deles conclusivo.

\end{enumerate}

\begin{destaque}[Cinco das seis limitações são do conjunto de dados, não do método]
Isso é característico da área. A comunidade de detecção de anomalias industriais dispõe
de poucos conjuntos públicos, e o 3W --- apesar de todas as ressalvas --- é um dos
melhores existentes.

Avançar exige mais dados públicos de qualidade, com anotação por especialista e taxonomia
por mecanismo. É trabalho menos glamouroso que propor arquiteturas, e provavelmente mais
impactante.
\end{destaque}

\section{Limitações de método}

\subsection{Segmentação}

A U-Net exige anotação por amostra, muito mais cara que anotação por janela. O
\cref{cap:segmentacao} obteve $0{,}863$ de acurácia com esse custo. Não se sabe qual
seria o desempenho de abordagens fracamente supervisionadas, que aprendem a segmentar a
partir de rótulos de janela.

A revocação SoftED de $0{,}678$ significa que cerca de um terço dos eventos não é
detectado dentro da tolerância temporal. A detecção antecipada de $0{,}222$ indica que o
modelo raramente antecipa o especialista --- e antecipar é o que teria valor operacional.

\subsection{Novidade}

A média $\bm{\mu}$ e a covariância $\bm{\Sigma}$ do Mahalanobis são estimadas uma vez, e
não acompanham deriva. Um esquema de atualização incremental é necessário, mas introduz o
risco de absorver degradação progressiva --- o modo de falha silencioso do
\cref{cap:operacao}.

O limiar no percentil 95 é uma escolha, não uma derivação. A justificativa dada no
\cref{cap:mahalanobis} --- que a hipótese de normalidade multivariada não se sustenta ---
explica por que não se usa $\chi^2$, mas não estabelece que 95 seja ótimo.

A classe 7 permanece em $52{,}3\,\%$. O diagnóstico é claro: a janela de duzentos passos
não comporta um fenômeno de semanas. A solução exige representação multiescala.

\subsection{Agente}

\begin{itemize}[leftmargin=*,itemsep=0.3em]
  \item \textbf{Sem ablação.} Não se sabe quais das treze métricas importam. Um estudo de
        ablação é o experimento de menor custo e maior retorno informativo pendente.
  \item \textbf{Sem comparação com atribuição.} SHAP e gradientes integrados produzem
        explicações de natureza diferente; não foram comparados.
  \item \textbf{Sem ajuste fino, sem exemplos, sem múltiplas passagens, sem ferramentas.}
        Como o \cref{cap:agente} registrou, os resultados devem ser lidos como piso.
  \item \textbf{Sem calibração.} A confiança declarada pelo agente não foi verificada
        contra frequências empíricas.
  \item \textbf{Sem informação de localização.} A confusão entre as classes 8 e 9 decorre
        disso, e é corrigível.
\end{itemize}

\subsection{Reprodutibilidade do modelo de linguagem}

O \cref{cap:llm} prometeu que este capítulo discutiria alternativas. A questão é séria:
o modelo é servido por interface de programação de terceiros, pode ser atualizado ou
descontinuado sem aviso, e a temperatura de $0{,}2$ reduz mas não elimina a
estocasticidade. Os resultados não são estritamente reproduzíveis.

\begin{table}[htbp]
\centering
\caption{Alternativas para reprodutibilidade da camada de linguagem.}
\label{tab:fut-repro}
\small
\begin{tabular}{p{3.0cm}p{4.4cm}p{4.6cm}}
\toprule
\textbf{Alternativa} & \textbf{Ganho} & \textbf{Custo} \\
\midrule
Modelo aberto local, versão fixa &
Reprodutibilidade estrita; dados não saem do perímetro &
Infraestrutura de GPU; desempenho tipicamente inferior \\
\addlinespace
Ajuste fino sobre modelo aberto menor &
Adaptação ao domínio; modelo próprio e versionável &
Exige dados rotulados --- o que o sistema procura evitar \\
\addlinespace
Registro completo de entradas e saídas &
Auditoria posterior; permite reanálise &
Não permite repetir o experimento \\
\addlinespace
Amostragem múltipla com agregação &
Reduz variância; mede consistência &
Multiplica o custo por chamada \\
\bottomrule
\end{tabular}
\end{table}

\begin{destaque}[Não há escolha dominante]
Cada linha da \cref{tab:fut-repro} troca uma propriedade desejável por outra. A
recomendação prática é combinar a terceira linha --- registro completo, que é barata e
compatível com qualquer das outras --- com a alternativa que o contexto de governança
exigir.

Em ambiente industrial com restrição de dados, a primeira linha tende a ser mandatória,
e o custo de desempenho deve ser aceito.
\end{destaque}

\section{Limitações de avaliação}

\begin{enumerate}[leftmargin=*,itemsep=0.4em]

\item \textbf{Nenhuma avaliação com operadores.} Todas as afirmações sobre carga
cognitiva, fadiga de alarme e utilidade das justificativas são argumentadas, não medidas.
A política de triagem do \cref{cap:resultados-agente} --- que reduziria à metade o volume
de revisão --- é uma projeção.

\item \textbf{Nenhuma avaliação fim-a-fim.} Os componentes foram medidos isoladamente. A
arquitetura de referência do \cref{cap:sintese} nunca foi executada como um todo.

\item \textbf{Ausência de linha de base humana.} Não se sabe qual é a acurácia de um
especialista sobre os mesmos segmentos, o que impede afirmar que o sistema é útil e não
apenas mensurável.

\item \textbf{Comparações com a literatura são aproximadas.} Como o \cref{cap:binarios}
registrou, protocolos, partições e definições de classe diferem entre trabalhos.

\end{enumerate}

\begin{destaque}[A limitação mais séria é a terceira]
Sem linha de base humana, ``$89{,}7\,\%$ de detecção de novidade'' é um número sem
referencial. Pode ser excelente ou trivial.

Estabelecer essa linha de base é caro --- exige tempo de especialista --- e é a lacuna
que mais compromete a interpretação de todo o trabalho.
\end{destaque}

\section{Uma agenda}

\subsection{Ordenada por razão informação/custo}

\begin{table}[htbp]
\centering
\caption{Agenda de trabalho, ordenada por razão entre informação obtida e custo.}
\label{tab:fut-agenda}
\small
\begin{tabular}{p{5.2cm}p{2.0cm}p{4.8cm}}
\toprule
\textbf{Experimento} & \textbf{Custo} & \textbf{Informação obtida} \\
\midrule
Remoção do confundimento de telemetria & Baixo &
Quantifica quanto do desempenho é artefato \\
\addlinespace
Ablação das treze métricas & Baixo &
Quais métricas importam; guia para adicionar novas \\
\addlinespace
Métricas de localização por sensor & Baixo &
Testa a hipótese sobre a confusão entre classes 8 e 9 \\
\addlinespace
Calibração da confiança declarada & Baixo &
Viabiliza uso da confiança na política de triagem \\
\addlinespace
Múltiplas passagens com agregação & Médio &
Mede consistência; provável ganho em top-1 \\
\addlinespace
Representação multiescala & Médio &
Ataca a classe 7 diretamente \\
\addlinespace
Avaliação fim-a-fim da arquitetura & Médio &
Único número que descreve o sistema real \\
\addlinespace
Segmentação fracamente supervisionada & Alto &
Elimina o custo de anotação por amostra \\
\addlinespace
Linha de base humana & Alto &
Referencial para todos os números do livro \\
\addlinespace
Redefinição das classes sintomáticas & Alto &
Corrige a taxonomia; beneficia toda a comunidade \\
\bottomrule
\end{tabular}
\end{table}

\subsection{Quatro questões de pesquisa}

\begin{description}[leftmargin=0pt,style=nextline,itemsep=0.5em]

\item[A assimetria é universal?] Métodos geométricos e modelos de linguagem, sistemas de
natureza completamente distinta, exibem a mesma assimetria entre afirmar identidade e
reconhecer incompatibilidade. Ela se manifesta em outros domínios? Existe uma formulação
teórica que a preveja?

\item[A convergência de nomes é uma métrica de confiança válida?] O
\cref{cap:resultados-agente} propôs medir a concentração das descrições geradas como
indicador não supervisionado de confiabilidade. A proposta é plausível e não foi
validada. Se funcionar, resolve o problema de estimar confiança sem rótulos.

\item[Quanto da vantagem do latente supervisionado é transferível?] O
\cref{cap:mahalanobis} mostrou que representar com um classificador em vez de um
autoencoder eleva a detecção de novidade em quarenta e seis pontos. Isso depende de haver
classes conhecidas suficientes. Qual é o mínimo? A vantagem persiste com três classes?
Com duas?

\item[Sistemas que descrevem auditam melhor os dados?] O artefato de telemetria do 3W só
se tornou visível pela descrição em linguagem natural. Isso sugere um uso não previsto de
modelos de linguagem: auditoria automática de conjuntos de dados. A hipótese pode ser
testada injetando artefatos conhecidos e verificando se são descritos.

\end{description}

\section{Encerramento}

\begin{destaque}[O que este livro sustenta]
Um sistema de monitoramento pode --- e deve --- dizer ``não sei''.

Dizer isso corretamente é mais fácil do que dizer o que algo é: $95{,}4\,\%$ contra
$87\,\%$ nos métodos numéricos, $89{,}7\,\%$ contra $35{,}1\,\%$ na camada de linguagem.

O espaço em que a pergunta é feita importa mais que o método: a mesma distância detecta
$67{,}5\,\%$ ou $98{,}2\,\%$ da mesma anomalia.

E um sistema que descreve o que vê, em linguagem que o especialista lê, transforma um
grupo anônimo em uma hipótese verificável --- além de revelar problemas nos dados que
nenhuma métrica agregada exporia.

Nada disso substitui o especialista. Tudo isso muda o que ele recebe no início do turno.
\end{destaque}

\resumocapitulo{%
Os resultados deste livro devem ser lidos sob uma ressalva que atravessa todos os
capítulos: no 3W, o início de muitas anomalias coincide com falha de telemetria, e parte
do desempenho reportado pode decorrer da detecção desse artefato --- fração que os
experimentos realizados não permitem quantificar. Cinco das seis limitações de dados são
do conjunto, não do método, o que aponta a escassez de dados públicos de qualidade como
gargalo da área. Entre as limitações de método, destacam-se a revocação SoftED de
$0{,}678$ e a detecção antecipada de $0{,}222$ na segmentação; a estimação estática de
$\bm{\mu}$ e $\bm{\Sigma}$ no Mahalanobis; a classe 7 presa em $52{,}3\,\%$ por
insuficiência de escala temporal; e, no agente, a ausência de ablação, de calibração e de
métricas de localização. A reprodutibilidade da camada de linguagem não tem solução
dominante: cada alternativa troca uma propriedade desejável por outra, e o registro
completo de entradas e saídas é a única compatível com todas as demais. A limitação de
avaliação mais séria é a ausência de linha de base humana, sem a qual os números do livro
carecem de referencial. A agenda proposta ordena dez experimentos por razão entre
informação e custo, com a avaliação do possível confundimento de telemetria em primeiro
lugar, e formula quatro
questões de pesquisa --- entre elas, se a assimetria entre reconhecer incompatibilidade e
afirmar identidade é uma propriedade universal.}

\begin{projetospesquisa}
\item Projete e execute o estudo de ablação das treze métricas do \cref{cap:agente}.
      Ordene-as por contribuição marginal. O resultado é consistente com a intuição
      física?

\item Implemente o protocolo de remoção do confundimento de telemetria e reavalie os
      resultados dos \cref{cap:binarios,cap:mahalanobis,cap:resultados-agente}. Quanta
      acurácia se perde em cada um? A ordenação entre métodos se mantém?

\item Verifique a calibração da confiança declarada pelo agente construindo um diagrama
      de confiabilidade. Se ela estiver mal calibrada, proponha e avalie uma correção.

\item Ataque a classe 7 com uma representação multiescala --- por exemplo, concatenando
      latentes de janelas de duzentos, dois mil e vinte mil passos. Quanto se ganha? O
      que se perde nas demais classes?

\item Estime experimentalmente o número mínimo de classes conhecidas necessário para que
      o latente supervisionado supere o latente do autoencoder na detecção de novidade.

\item Teste a hipótese de auditoria automática: injete um artefato conhecido em um
      conjunto de dados limpo e verifique se o agente o descreve. Varie a intensidade do
      artefato e determine o limiar de detecção.

\item Projete o protocolo de uma avaliação com operadores que meça, e não apenas
      argumente, a redução de carga de revisão prometida pela política de triagem do
      \cref{cap:resultados-agente}. Que ameaças à validade ele precisa controlar?
\end{projetospesquisa}

\begin{leituras}
\item \citet{Pimentel2014review} --- levantamento que contextualiza as limitações
      discutidas neste capítulo.
\item \citet{Salles2024softed} --- a métrica cujos valores de revocação e detecção
      antecipada apontam as lacunas da segmentação.
\item \citet{VazVargas2026_3w2} --- a evolução do 3W, e o caminho pelo qual várias das
      limitações de dados podem ser atacadas.
\item \citet{Sharma2025xaiagents}, \citet{Belay2026agenticAD} e \citet{Zhang2026llmrules}
      --- direções contemporâneas para a camada de linguagem.
\end{leituras}
